\chapter{What's New in Bijou 7.2.0}\label{record-283}

\noindent{\footnotesize\color{BookMuted}\path{docs/releases/7.2.0/whats-new.md}\par}\smallskip

Bijou 7.2.0 is a stabilization release for the V7 demo and documentation surface. It keeps the public Runtime Graph direction intentionally narrow while tightening framework input behavior, DOGFOOD readability, Blocks documentation, release evidence, and dependency review.

The public tag was created from the verified 2026-07-05 \texttt{m\allowbreak{}a\allowbreak{}i\allowbreak{}n} release commit. This document is the versioned release overview DOGFOOD and the repository front door use for 7.2.0.

\section{Framework Input Stabilization}\label{record-283-framework-input-stabilization}

The TUI framework now has better coverage for page-scoped input fallthrough and scripted mouse sequences. Workspace mouse movement, release, wheel, and non-left press events can be replayed through deterministic driver helpers, and the public TUI package exposes the frame helper APIs needed by application code.

This is the practical fix for release-demo workflows that need reliable mouse inspection without relying on terminal-specific manual behavior.

\section{DOGFOOD Readability And Release Story}\label{record-283-dogfood-readability-and-release-story}

DOGFOOD received another pass over the surfaces that made the 7.1 release-video rehearsal hard to inspect:

\begin{itemize}
\tightlist
\item
  locale fallback copy now distinguishes missing translations from intentional English source text
\item
  light-theme screenshots and fixture coverage keep color contrast and shell state readable
\item
  in-app release-story pages now point reviewers at the V7.2 evidence packet rather than asking them to infer the release state from scattered docs
\end{itemize}

The goal is not a new product surface. The goal is that the existing proof surface tells the truth when somebody records or reviews it.

\section{Blocks Documentation}\label{record-283-blocks-documentation}

The Blocks app-binding guide now includes copyable snippets and deterministic fixtures for the current binding shape. The documentation remains scoped to the V7 proof path, but it gives readers enough structure to understand how bindings, actions, and DOGFOOD examples relate before the larger V8 work starts.

\section{Theme Variant Coverage}\label{record-283-theme-variant-coverage}

First-party dark and light theme variants have explicit coverage in the release packet. The release keeps theme behavior focused on readability and replay evidence rather than opening a broader design-system expansion.

\section{Release Evidence}\label{record-283-release-evidence}

The release gate now carries a versioned evidence packet alongside the package metadata bump. That packet records trackers, landed PRs, replay commands, security audit scope, registry verification, release-readiness posture, and residual risk.

Maintainers should treat the release evidence packet as the authoritative 7.2 release record: it includes the public GitHub Release, final-main validation, publish workflow, tag CI, Tag Guard, and npm registry verification evidence.

\section{V8 Posture}\label{record-283-v8-posture}

Bijou 7.2.0 intentionally stops short of the next Runtime Graph product. The forward path remains:

% Presentation of original arrow notation. See flow-manifest.json.
\begin{center}
\begin{tikzpicture}[book flow node/.style={draw=BookHeading,line width=.65pt,fill=BookPanel,rounded corners=1.5pt,align=center,inner sep=5pt,minimum height=10mm,font=\fontsize{9}{11}\selectfont,execute at begin node={\hyphenpenalty=10000\exhyphenpenalty=10000}},book flow edge/.style={-{Latex[length=2mm]},draw=BookHeading,line width=.7pt}]
\node[book flow node,text width=58mm] (n0) at (0,0.0) {Graph\allowbreak{}QL Blocks};
\node[book flow node,text width=58mm] (n1) at (0,-1.55) {Bijou Blocks IR};
\node[book flow node,text width=58mm] (n2) at (0,-3.1) {Geordi render targets};
\node[book flow node,text width=58mm] (n3) at (4,-3.1) {TUI};
\node[book flow node,text width=58mm] (n4) at (4,-4.9) {Browser};
\node[book flow node,text width=58mm] (n5) at (4,-6.7) {Native UI};
\draw[book flow edge] (n0) -- (n1);
\draw[book flow edge] (n1) -- (n2);
\draw[book flow edge] (n2) -- (n3);
\draw[book flow edge] (n2) -- (n4);
\draw[book flow edge] (n2) -- (n5);
\end{tikzpicture}
\end{center}


That work belongs to V8 and later. V7.2 makes the current release surface clean enough to move on without carrying avoidable release debt forward.

